\section{引言：什么是 Harness Engineering}

LangChain 团队在这篇博客中分享了一个重要发现：在不更换底层模型（GPT-5.2-Codex）的前提下，仅通过改进 Agent Harness（即围绕模型构建的系统层——包括系统提示词、工具选择和执行流程），就将其编程 Agent（deepagents-cli）在 Terminal Bench 2.0 基准测试上的得分从 52.8\% 提升到 66.5\%，排名从 Top 30 跃升至 Top 5。

这一结果的核心启示是：\textbf{模型能力的释放程度很大程度上取决于 Harness 的设计质量}。

\begin{importantbox}{Harness Engineering 的核心定义}
Harness 的目标是将模型本身"参差不齐的智能"（spiky intelligence）塑造为适配目标任务的稳定能力。Harness Engineering 关注的是围绕模型构建的\textbf{系统工程}，包括系统提示词（System Prompt）、工具集合（Tools）、执行流（Execution Flow）等。优化目标包括任务完成率、Token 效率、延迟等。
\end{importantbox}

\subsection{为什么 Harness 如此重要}

当前的大语言模型本质上是黑箱：我们难以直接理解其内部推理机制。但 Harness 设计者可以控制的是模型的\textbf{输入}（Prompt、上下文注入）和对模型\textbf{输出}的后处理（中间件、循环检测等）。通过精心设计这些外部机制，可以在不改变模型权重的情况下显著提升实际表现。

LangChain 的方法论很简单：
\begin{enumerate}[leftmargin=2em]
    \item 用 Traces（LangSmith 中的执行追踪）理解 Agent 的失败模式
    \item 针对失败模式，调整 Harness 的三个旋钮：System Prompt、Tools、Middleware
    \item 迭代验证改进效果
\end{enumerate}

\begin{knowledgebox}{Terminal Bench 2.0 简介}
Terminal Bench 2.0 是当前评估 Agentic Coding 能力的标准基准测试，包含 89 个任务，涵盖机器学习、调试、生物信息学等领域。测试流程由 Harbor 编排：为每个任务启动沙盒环境（Daytona），与 Agent 交互，运行验证和评分脚本。所有 Agent 操作均存储在 LangSmith 中，包括延迟、Token 计数和成本等指标。
\end{knowledgebox}

\subsection{本章小结}

Harness Engineering 的核心理念是：模型智能是"原材料"，Harness 是"加工工艺"。同一个模型在不同 Harness 下的表现可以天差地别。LangChain 通过 Trace 驱动的迭代改进，展示了一套系统化的 Harness 优化方法论。


